% GENERATED by scripts/build-lecture-slides.py; do not edit. Portable profile: XeTeX, theme and fonts from the core.
\makeatletter\IfFileExists{core/beamer/beamerthememx.sty}{\def\mxcorepath{core/}\def\input@path{{core/beamer/}{core/}}}{\def\mxcorepath{../core/}\def\input@path{{../core/beamer/}{../core/}}}\makeatother
\documentclass[aspectratio=169,10pt,t]{beamer}
\usetheme{mx}
\usepackage{graphicx,booktabs}
\usepackage[french]{babel}
\setlength{\parskip}{.45em}
\hypersetup{pdfauthor={},pdfsubject={Lecture slides},pdfkeywords={},pdftitle={Similitudes planes}}
\providecommand{\N}{\mathbb{N}}
\providecommand{\Z}{\mathbb{Z}}
\providecommand{\D}{\mathbb{D}}
\providecommand{\Q}{\mathbb{Q}}
\providecommand{\R}{\mathbb{R}}
\providecommand{\C}{\mathbb{C}}
\providecommand{\vect}[1]{\overrightarrow{#1}}
\providecommand{\vi}{\vec{\imath}}
\providecommand{\vj}{\vec{\jmath}}
\providecommand{\vk}{\vec{k}}
\providecommand{\norm}[1]{\left\lVert #1\right\rVert}
\providecommand{\vc}[2]{\begin{pmatrix}#1\\#2\end{pmatrix}}
\providecommand{\deter}[2]{\det\left(#1,#2\right)}
\providecommand{\abs}[1]{\left\lvert #1\right\rvert}
\providecommand{\dd}{\,\mathrm{d}}
\providecommand{\e}{\mathrm{e}}
\providecommand{\Card}{\operatorname{Card}}
\providecommand{\E}{\operatorname{E}}
\providecommand{\V}{\operatorname{V}}
\providecommand{\degre}{^{\circ}}
\providecommand{\un}[1]{\,\mathrm{#1}}
\providecommand{\sci}[2]{#1\times10^{#2}}
\providecommand{\pgcd}{\operatorname{PGCD}}
\providecommand{\ppcm}{\operatorname{PPCM}}
\title{Similitudes planes}
\subtitle{Cours et exemples}
\date{}
\author{}
\begin{document}
\mxtitleframe
\begin{frame}{Généralités}
Une similitude plane \(S\) est une transformation géométrique du plan pour laquelle il existe un réel \(k>0\) tel que,
            pour tous points \(M\) et \(N\) du plan, en posant \(M' = S(M)\) et \(N' = S(N)\) :
                  \[
                        M'N' = k\, MN.
                  \]

Le réel \(k\) est appelé le \emph{rapport} de la similitude \(S\).

Toute isométrie est une similitude de rapport \(1\).
            Une homothétie de rapport \(k\) (\(k \neq 0\)) est une similitude de rapport \(|k|\).
\end{frame}
\begin{frame}{Généralités}
Toute similitude \(S\) de rapport \(k (>0)\) peut s'écrire sous la forme \(S = h \circ f\)
                  où \(h\) est une homothétie de rapport \(k\) et \(f\) une isométrie.

\begin{itemize}

                        \item Si \(f\) est un déplacement alors \(S\) est une similitude plane directe.

\item Si \(f\) est un antidéplacement alors \(S\) est une similitude plane indirecte.

\end{itemize}
\end{frame}
\begin{frame}{Généralités}
Soit \(S\) une similitude de rapport \(k\).
                        Soit \(h\) l'homothétie de centre \(I\) et de rapport \(k\).
                        On a \(h^{-1} = h\left(I ; \dfrac{1}{k}\right)\).
                        On pose \(f = h^{-1} \circ S\).
                        
                        Montrons que \(f\) est une isométrie : 
                              Soient \(M\) et \(N\) deux points, posons \(M' = f(M)\) et \(N' = f(N)\).
                              On a alors 
                              \[\begin{aligned}
                                    M' & = h^{-1} \left[ S(M) \right]       & N' & = h^{-1} \left[ S(N) \right]      \\
                                       & = h^{-1} (M_1)                     &    & = h^{-1} (N_1)                    \\
                              \end{aligned}\]
                              où \(M_1 = S(M)\) et \(N_1 = S(N)\), et donc \(M_1 N_1 = k MN\).
                              Ainsi, \(M'N' = \dfrac{1}{k} M_1 N_1 = \dfrac{1}{\not k} \not k MN\).
                              On déduit que \(M'N' = MN\), et par conséquent, que \(f\) est une isométrie.
                              
                        Montrons que \(S = h\circ f\) : 
                              \[
                                    f = h^{-1} \circ S 
                                          \Rightarrow 
                                                h \circ f = \underbrace{h \circ h^{-1}}_{\mathrm{Id}} \circ S
                                                      \Rightarrow
                                                            h \circ f = S.
                              \]
\end{frame}
\begin{frame}{Généralités}
\begin{enumerate}
\setcounter{enumi}{0}\item Toute similitude est une application affine bijective.

\end{enumerate}
\end{frame}
\begin{frame}{Généralités (2)}
\begin{enumerate}
\setcounter{enumi}{1}\item La composition de \(2\) similitudes est une similitude.
                              Si
                              
\begin{itemize}

                                    \item \(S\) est une similitude de rapport \(k\),
                                    
\item \(S'\) une similitude de rapport \(k'\),

\end{itemize}

                              alors \(S \circ S'\) est une similitude de rapport \(kk'\).

\end{enumerate}
\end{frame}
\begin{frame}{Généralités (3)}
\begin{enumerate}
\setcounter{enumi}{4}\item Une similitude de rapport \(k\) multiplie les distances par \(k\) et les aires par \(k^2\).
                        
\item Une similitude conserve le parallélisme, l'orthogonalité, l'alignement et les milieux ;
                              une similitude directe conserve les angles orientés, une similitude indirecte les change en leurs opposés.

\end{enumerate}
\end{frame}
\begin{frame}{Similitude plane directe}
Une similitude plane directe \(S\) est soit une translation
            soit la composition d'une homothétie et d'une rotation de même centre :
            c'est la forme réduite.

Une similitude plane directe \(S\) a pour expression complexe \(z' = az + b\) avec \(a \in \mathbb{C}^*\) et \(b \in \mathbb{C}\).
\end{frame}
\begin{frame}{Similitude plane directe (2)}
\begin{itemize}

                  \item \(|a| = k\) est le rapport,
                  
\item \(\mathrm{Arg}\left( a \right) = \alpha\) est l'angle,
                  
\item \(I\left(\dfrac{b}{1-a}\right)\) est le centre (si \(a\neq 1\)).

\end{itemize}
\end{frame}
\begin{frame}{Similitude plane directe (3)}
Une similitude plane directe est définie analytiquement par un système du type :
	            \[
	                  \left\lbrace
	                  \begin{array}{l}
	                        x' = a x - b y + x_0, \\
	                        y' = b x + a y + y_0.
	                  \end{array}
	                  \right.
	            \]

où \(a^2 + b^2 = k^2\), avec \(k > 0\) le rapport.
\end{frame}
\begin{frame}{Similitude plane indirecte}
Une similitude plane indirecte \(f\) est la composée d'une homothétie et d'un antidéplacement.
      Le rapport \(k\) de \(f\) est celui de l'homothétie (en valeur absolue).
\end{frame}
\begin{frame}{Similitude plane indirecte (2)}
\begin{itemize}
\item Si \(k=1\) :  \(f\) est un antidéplacement.
                  L'ensemble des points invariants de \(f\) est soit vide soit une droite.
                  
\begin{itemize}

                        \item Si \(f\) n'a pas de point invariant, alors \(f\) est la composée d'une symétrie
                              orthogonale et d'une translation de vecteur non nul.
                        
\item Si l'ensemble des points invariants est une droite \((\mathcal{D})\), alors 
                              \(f\) la symétrie orthogonale d'axe \((\mathcal{D})\).

\end{itemize}

\end{itemize}
\end{frame}
\begin{frame}{Similitude plane indirecte (3)}
\begin{itemize}
\item Si \(k\neq 1\) :  \(f\) admet un seul point invariant \(\Omega\).
                  Il existe une droite \((\mathcal{D})\) passant par \(\Omega\) telle que 
                  \(f = h(\Omega,k) \circ S_{(\mathcal{D})} = S_{(\mathcal{D})} \circ h(\Omega,k)\).
                  
                  La droite \((\mathcal{D})\) est l'ensemble des points \(M\) tels que 
                  \(\overrightarrow{\Omega M'} = k \cdot \overrightarrow{\Omega M}\),

                  c'est la \emph{première droite globalement invariante} par \(f\).
                  
                  La droite \((\mathcal{D}') = \{ M \mid \overrightarrow{\Omega M'} = -k \cdot \overrightarrow{\Omega M} \}\)

                  est la \emph{deuxième droite globalement invariante} par \(f\).

\end{itemize}
\end{frame}
\begin{frame}{Similitude plane indirecte}
Soit \(f\) une similitude plane indirecte de rapport \(k\).
            La similitude \(f\) est alors définie analytiquement par un système du type :
            \(\left\lbrace \begin{array}{l} x' = a x + b y + x_0 \\ y' = b x - a y + y_0 \end{array} \right.\)
            où \(a^2 + b^2 = k^2\).
\end{frame}
\begin{frame}{Similitude plane indirecte}
Soit \(S\) la similitude plane indirecte égale à \(h \circ S_{(\mathcal{D})}\), où \(h\) est une homothétie de rapport \(k\).
            Soient 
            \[
                  \begin{array}{l}
                        h : z' = k z + b \quad (k \in \mathbb{R}^*) \\
                        S_{(\mathcal{D})} : z' = a \overline{z} + c 
                  \end{array}
            \]

les expressions complexes de \(h\) et \(S_{(\mathcal{D})}\).

On a alors, \(S : z' = k \left( a \overline{z} + c \right) + b = k a \overline{z} + k c + b\).
\end{frame}
\begin{frame}{Similitude plane indirecte (2)}
De plus, on vérifie que \(|k a| = |k| \cdot |a| = |k|\) (\(S_{(\mathcal{D})}\) étant une isométrie, \(|a| = 1\)).

Une similitude plane indirecte a une expression complexe de la forme \(z' = a \overline{z} + b\) (\(a \neq 0\))

et son rapport est \(k = |a|\).
\end{frame}
\begin{frame}{Similitude plane indirecte}
Le point invariant \(I\)

 Toute similitude plane indirecte de rapport \(k \neq 1\) admet un point invariant \(I\) qu'on appelle le \emph{centre}.
\end{frame}
\begin{frame}{Similitude plane indirecte (2)}
\[
                        \left.
                              \begin{array}{l}
                                    z_I = a \overline{z_I} + b \\
                                    \overline{z_I} = \overline{a} z_I + \overline{b} 
                              \end{array}
                        \right\rbrace
                              \Rightarrow 
                        z_I = a (\overline{a} z_I + \overline{b} ) + b.
                  \]
\end{frame}
\begin{frame}{Similitude plane indirecte (3)}
Donc \(z_I = a \overline{a} z_I + a \overline{b} + b\).
            Il vient que \((1-a\overline{a}) z_I = a\overline{b} + b\) et donc que

\[
                        z_I = \dfrac{a\overline{b}+b}{1 - a \overline{a}} (|a| \neq 1)
                  \]

est l'affixe du centre \(I\).
\end{frame}
\begin{frame}{Similitude plane indirecte}
Soit \(S\) la transformation géométrique d'expression complexe \(z'=2\overline{z}+3i\).
                  Quelle est la nature de \(S\) ? Déterminer ses éléments caractéristiques.
                  
                  Construire ensuite le transformé du cercle \((\mathcal{C})\) de centre \(A(1-i)\) et de rayon 1 par \(S\).
                  Ecrire l'équation de ce cercle.
\end{frame}
\end{document}
